\section{Claims and gates}
Study~0 maintains a claim ladder rather than one ''success'' verdict.

\begin{description}[leftmargin=2.2cm,style=nextline]
\item[C0 Feasibility] Can contracts be elicited/resolved with low burden and low interference?
\item[C1 Forecast validity] Are the resulting probabilities measurable with appropriate probabilistic diagnostics?
\item[C2 Incremental aggregation] What is the paired market-vs-mean Brier difference?
\item[C3 Information condition] Exploratorily, when does market aggregation appear most useful?
\item[C4 Human--machine complementarity] Future work only.
\end{description}

Operational-use consideration is blocked by design until separate evidence and review authorize a later intervention study.

\section{Stopping and prospective freeze}
A deployment must freeze either a calendar window or a target number of resolved contracts plus a maximum calendar end. Outcome-dependent early stopping is prohibited. Approximately 20--40 resolved contracts is a pragmatic feasibility target, not a power guarantee. If fewer than 20 admissible contracts resolve, the study reports feasibility/descriptive evidence and does not promote C2 as established.

The following must be frozen before outcome inspection: contract semantics, participant eligibility, market configuration, baselines, primary estimand, scoring rule, void/interference rule, stopping rule and random resampling seed.

\section{Organizational safeguards}
The experiment concerns employees and operational knowledge, so organizational validity is inseparable from scientific validity. Required controls include voluntary participation, no individual HR use, no real-money betting by default, pseudonymous research identifiers, data minimization, non-retaliation for pessimistic forecasts, objective resolution and a stop condition if market outputs begin to alter live decisions.

Raw internal rationales can contain confidential architecture or personnel information and are therefore excluded by default. The public repository contains only generic examples and schemas. Any live deployment must use an approved private data store and normal organizational privacy/security/compliance processes.

\section{Threats to validity}
\paragraph{Construct validity.} A market price is an aggregate forecast under a mechanism, not truth. A perfectly resolvable question may still be irrelevant or gameable.

\paragraph{Internal validity.} Forecast exposure can change behavior; insiders may already know outcomes; question authors can leak expectations; resolution discretion can introduce post hoc bias.

\paragraph{Statistical validity.} Study~0 is small, contracts can be dependent, rare outcomes can make calibration unstable, and exploratory analyses are multiplicity-prone.

\paragraph{External validity.} Corporate-market and new-product results do not automatically generalize to engineering qualification, and one organization, product family or market platform may not generalize to another.

\paragraph{Platform validity.} Thin participation, market-maker settings, hidden product updates and AI/news assistance can change the measured signal.

\paragraph{Causal validity.} Conditional forecasts $P(Y\mid A)$ are not automatically causal interventions $P(Y\mid do(A))$. Study~0 therefore does not use conditional markets to rank engineering actions.

\section{Instrument candidate: mature internal-market platform}
Building a bespoke market engine would conflate two questions: whether collective forecasting adds information and whether a new software prototype works. A mature enterprise platform can be a cleaner Study~0 instrument if security, identity, export and algorithmic reproducibility requirements are met. A platform such as Hypermind/Prescience may therefore be evaluated as an instrument, but the protocol remains vendor-neutral and vendor accuracy claims are not scientific evidence for this context.

Preferred Study~0 configuration disables automated external news/web/AI assistance, uses virtual points, captures private forecasts before market exposure, withholds aggregate outputs from decision owners, minimizes free text and exports the terminal market state plus configuration immutably.

\section{Natural research connections without forced coupling}
\subsection{Executable/evidence contracts}
The Forecast Contract naturally complements evidence-centric engineering: a forecast resolves against an authoritative evidence object and frozen semantics. This is a conceptual compatibility, not a dependency on a particular executable-contract implementation.

\subsection{Probabilistic engineering models}
If a qualified model already estimates an engineering outcome from structured signals, later work can compare
\[
 p^{human},\quad p^{market},\quad p^{model},\quad p^{hybrid}.
\]
The scientifically interesting question is not ''human or AI?'' but whether errors are complementary and whether a combination rule generalizes to held-out contracts.

\subsection{Sequential evidence acquisition}
A future POMDP-like formulation could treat a collective forecast as an observation or treat eliciting one as an information-acquisition action with cost and expected value of information. This is a later hypothesis. GO-ED-POMDP is intentionally absent from the Study~0 estimand and causal design.

\section{Pedagogical contract}
The work also externalizes difficult concepts for reuse: prediction market, proper scoring rule, Brier score, calibration/sharpness and Forecast Contract. Each concept should be teachable through plain-language intuition, a concrete example, mathematical descent, an executable check, a misconception and an understanding gate. Canonical teaching versions belong in the Diderot MMALS/ML knowledge repository under independent source/provenance review; they do not become scientific authority merely by being pedagogically clear.

\section{Research roadmap}
\textbf{Study~0} measures feasibility and incremental forecast information in shadow mode. \textbf{Study~1} should replicate across event families, test robustness and predeclare actor/observer or diversity hypotheses. \textbf{Study~2} may evaluate human--machine ensembles using held-out contracts. \textbf{Study~3} may finally study whether exposing forecasts improves evidence acquisition or decisions under an intervention design that explicitly models behavioral feedback.

\section{Conclusion}
Prediction markets are attractive because they offer a compact mechanism for turning dispersed beliefs into a probabilistic signal. That attractiveness is precisely why an engineering organization should qualify the mechanism before operationalizing it. The proposed Study~0 treats collective forecasting as a measurement problem: define the event, preserve independent baselines, freeze the mechanism, resolve from authoritative evidence, score at the correct unit of analysis, record interference and preserve null/negative results. The desired output is not a new gate. It is evidence about whether a new evidence sensor deserves a future role.

\begin{thebibliography}{99}
\bibitem[Brier(1950)]{brier1950} Brier, G. W. (1950). Verification of Forecasts Expressed in Terms of Probability. \emph{Monthly Weather Review}, 78(1), 1--3.
\bibitem[Chen and Plott(2002)]{chenplott2002} Chen, K.-Y. and Plott, C. R. (2002). Information Aggregation Mechanisms: Concept, Design and Implementation for a Sales Forecasting Problem. California Institute of Technology; field deployment at Hewlett-Packard.
\bibitem[Hanson(2007)]{hanson2007} Hanson, R. (2007). Logarithmic Market Scoring Rules for Modular Combinatorial Information Aggregation. \emph{Journal of Prediction Markets}, 1(1), 3--15. doi:10.5750/jpm.v1i1.417.
\bibitem[Gneiting and Raftery(2007)]{gneiting2007} Gneiting, T. and Raftery, A. E. (2007). Strictly Proper Scoring Rules, Prediction, and Estimation. \emph{Journal of the American Statistical Association}, 102(477), 359--378. doi:10.1198/016214506000001437.
\bibitem[Arrow et al.(2008)]{arrow2008} Arrow, K. J. et al. (2008). The Promise of Prediction Markets. \emph{Science}, 320(5878), 877--878. doi:10.1126/science.1157679.
\bibitem[Cowgill et al.(2009)]{cowgill2009} Cowgill, B., Wolfers, J. and Zitzewitz, E. (2009). Using Prediction Markets to Track Information Flows: Evidence from Google. \emph{Auctions, Market Mechanisms and Their Applications}.
\bibitem[Matzler et al.(2013)]{matzler2013} Matzler, K., Grabher, C., Huber, J. and Fueller, J. (2013). Predicting New Product Success with Prediction Markets in Online Communities. \emph{R\&D Management}, 43(5), 420--432. doi:10.1111/radm.12030.
\bibitem[Cowgill and Zitzewitz(2015)]{cowgill2015} Cowgill, B. and Zitzewitz, E. (2015). Corporate Prediction Markets: Evidence from Google, Ford, and Firm X. \emph{The Review of Economic Studies}, 82(4), 1309--1341. doi:10.1093/restud/rdv014.
\bibitem[Atanasov et al.(2017)]{atanasov2017} Atanasov, P. et al. (2017). Distilling the Wisdom of Crowds: Prediction Markets vs. Prediction Polls. \emph{Management Science}, 63(3), 691--706. doi:10.1287/mnsc.2015.2374.
\bibitem[Qiu et al.(2017)]{qiu2017} Qiu, L., Cheng, H. K. and Pu, J. (2017). Hidden Profiles in Corporate Prediction Markets: The Impact of Public Information Precision and Social Interactions. \emph{MIS Quarterly}, 41(4), 1249--1273. doi:10.25300/MISQ/2017/41.4.11.
\bibitem[Hong et al.(2021)]{hong2021} Hong, L., Lamberson, P. J. and Page, S. E. (2021). Hybrid Predictive Ensembles: Synergies Between Human and Computational Forecasts. \emph{Journal of Social Computing}, 2(2), 89--102. doi:10.23919/JSC.2021.0009.
\bibitem[Wang and Pfeiffer(2023)]{wang2023} Wang, W. and Pfeiffer, T. (2023). Proxy Forecasting to Avoid Stochastic Decision Rules in Decision Markets. arXiv:2303.10857.
\end{thebibliography}
